\ifdefined\gkver\else\def\gkver{student}\fi
\documentclass[\gkver]{gaokaozhenti}
\newcommand{\ccnum}[1]{{\fontspec{Noto Sans CJK JP}#1}}
\begin{document}

\chapter{2019年天津理综}
%% paperType: 理综物理部分

\begin{enumerate}
\item
%% number: 14
%% paperName: 2019年天津理综第 14 题 6 分
%% typeId: 1
%% type: 单选题
%% chapter: 曲线运动
%% point: 人造地球卫星
%% method: 无
%% score: 6
%% degree: 550
%% duplicateId: 0
%% body:
2018 年 12 月 8 日，肩负着亿万中华儿女探月飞天梦想的嫦娥四号探测器成功发射，“实现人类航天器首次在月球背面巡视探测，率先在月背刻上了中国足迹”。已知月球的质量为 $M$、半径为 $R$，探测器的质量为 $m$，引力常量为 $G$，嫦娥四号探测器围绕月球做半径为 $r$ 的匀速圆周运动时，探测器的\xzanswer{A}
\onepicture[h=4.2cm]{figs/14.jpg}
\fourchoices[answer=A]
{周期为 $\sqrt{\frac{4\pi^{2}r^{3}}{GM}}$}
{动能为 $\frac{GMm}{2R}$}
{角速度为 $\sqrt{\frac{Gm}{r^{3}}}$}
{向心加速度为 $\frac{GM}{R^{2}}$}
%% answer: A
%% memo:
\memoanswer{
根据万有引力提供卫星做匀速圆周运动的向心力，有 $\frac{GMm}{r^{2}}=ma=m\frac{v^{2}}{r}=m\omega^{2}r=m\frac{4\pi^{2}}{T^{2}}r$，整理得“嫦娥四号”绕月球做匀速圆周运动的周期 $T=\sqrt{\frac{4\pi^{2}r^{3}}{GM}}$，角速度 $\omega=\sqrt{\frac{GM}{r^{3}}}$，向心加速度 $a_{\text{向心}}=\frac{GM}{r^{2}}$，速度 $v=\sqrt{\frac{GM}{r}}$，故 A 正确 CD 错误；探测器的动能 $E_{k}=\frac{1}{2}mv^{2}=\frac{GMm}{2r}$，故 B 错误。
}

\item
%% number: 15
%% paperName: 2019年天津理综第 15 题 6 分
%% typeId: 1
%% type: 单选题
%% chapter: 力物体的平衡
%% point: 多力平衡
%% method: 无
%% score: 6
%% degree: 650
%% duplicateId: 0
%% body:
2018 年 10 月 23 日，港珠澳跨海大桥正式通车。为保持以往船行习惯，在航道处建造了单面索（所有钢索均处在同一竖直面内）斜拉桥，其索塔与钢索如图所示。下列说法正确的是\xzanswer{C}
\onepicture[h=4.6cm]{\includesvg{figs/15}}
\fourchoices[answer=C]
{增加钢索的数量可减小索塔受到的向下的压力}
{为了减小钢索承受的拉力，可以适当降低索塔的高度}
{索塔两侧钢索对称且拉力大小相同时，钢索对索塔的合力竖直向下}
{为了使索塔受到钢索的合力竖直向下，索塔两侧的钢索必须对称分布}
%% answer: C
%% memo:
\memoanswer{
对整座桥分析可知，索塔受到向下的压力等于桥的重力，与钢索的数量无关，故 A 错误；索塔越低，左右两侧钢索间的夹角越大，根据合力一定时，两对称分力间的夹角越大，则分力越大，可得钢索承受的力将变大，故 B 错误；如图所示，若 $F_1 = F_2$，且 $\theta_1 = \theta_2$，则由平行四边形定则可得合力方向一定竖直向下，故 C 正确；只要满足 $F_1 \sin \theta_1 = F_2 \sin \theta_2$，即水平分力相互抵消时，合力方向竖直向下，两侧的钢索不一定对称，故 D 错误。
\onepicture[h=3.0cm]{figs/15_an.jpg}
}

\item
%% number: 16
%% paperName: 2019年天津理综第 16 题 6 分
%% typeId: 1
%% type: 单选题
%% chapter: 电场
%% point: 带电粒子的偏转
%% method: 无
%% score: 6
%% degree: 550
%% duplicateId: 0
%% body:
如图所示，在水平向右的匀强电场中，质量为 $m$ 的带电小球，以初速度 $v$ 从 M 点竖直向上运动，通过 N 点时，速度大小为 $2v$，方向与电场方向相反，则小球从 M 运动到 N 的过程\xzanswer{B}
\onepicture[h=3.0cm]{figs/16.png}
\fourchoices[answer=B]
{动能增加 $\frac{1}{2}mv^{2}$}
{机械能增加 $2mv^{2}$}
{重力势能增加 $\frac{3}{2}mv^{2}$}
{电势能增加 $2mv^{2}$}
%% answer: B
%% memo:
\memoanswer{
带电小球受力分析如图所示，则水平方向上有 $qE = ma_x$①，$2v = a_x t$②，竖直方向上有 $mg = ma_y$，则 $a_y = g$③，$v - a_y t = 0$④，联立①②③④可得 $qE = 2mg$⑤，上升的高度 $h = \frac{v^2}{2g}$⑥，水平位移 $x = \frac{1}{2} \cdot 2v \cdot t = 2h$⑦，动能增加量 $\Delta E_k = \frac{1}{2} m(2v)^2 - \frac{1}{2} mv^2 = \frac{3}{2} mv^2$，故 A 错误；机械能的增加量 $\Delta E = qE \cdot x = 2mg \cdot 2\frac{v^2}{2g} = 2mv^2$，故 B 正确；重力势能增加量 $\Delta E_p = mgh = \frac{1}{2} mv^2$，故 C 错误；电场力做正功，电势能减小，则 $\Delta E_{\text{电}} = -qEx = -2mv^2$，即电势能减小 $2mv^2$，故 D 错误。
\onepicture[h=2.4cm]{figs/16_an.jpg}
}

\item
%% number: 17
%% paperName: 2019年天津理综第 17 题 6 分
%% typeId: 1
%% type: 单选题
%% chapter: 磁场
%% point: 洛伦兹力
%% method: 无
%% score: 6
%% degree: 550
%% duplicateId: 0
%% body:
笔记本电脑机身和显示屏对应部位分别有磁体和霍尔元件。当显示屏开启时磁体远离霍尔元件，电脑正常工作：当显示屏闭合时磁体靠近霍尔元件，屏幕熄灭，电脑进入休眠状态。如图所示，一块宽为 $a$、长为 $c$ 的矩形半导体霍尔元件，元件内的导电粒子是电荷量为 $e$ 的自由电子，通入方向向右的电流时，电子的定向移动速度为 $v$。当显示屏闭合时元件处于垂直于上表面、方向向下的匀强磁场中，于是元件的前、后表面间出现电压 $U$，以此控制屏幕的熄灭。则元件的\xzanswer{D}
\onepicture[h=3.0cm]{figs/17.png}
\fourchoices[answer=D]
{前表面的电势比后表面的低}
{前、后表面间的电压 $U$ 与 $v$ 无关}
{前、后表面间的电压 $U$ 与 $c$ 成正比}
{自由电子受到的洛伦兹力大小为 $\frac{eU}{a}$}
%% answer: D
%% memo:
\memoanswer{
由于导电粒子是自由电子，根据左手定则可知，电子向后表面偏转，则后表面得到电子带负电，电势较低，故 A 错误；霍尔效应达到平衡时洛伦兹力等于电场力，即 $evB = e\frac{U}{a}$，可得 $U = Bav$，故 BC 错误 D 正确。
}

\item
%% number: 18
%% paperName: 2019年天津理综第 18 题 6 分
%% typeId: 1
%% type: 单选题
%% chapter: 光学
%% point: 光电效应
%% method: 无
%% score: 6
%% degree: 450
%% duplicateId: 0
%% body:
如图为 a、b、c 三种光在同一光电效应装置中测的光电流和电压的关系。由 a、b、c 组成的复色光通过三棱镜时，下述光路图中正确的是\xzanswer{C}
\onepicture[h=3.0cm]{figs/18a.png}
\fourchoices[answer=C, ispicture=true, h=2.1cm]
{\includesvg{figs/18b}}
{\includesvg{figs/18c}}
{\includesvg{figs/18d}}
{\includesvg{figs/18e}}
%% answer: C
%% memo:
\memoanswer{
由光电效应方程得 $E_k = h\nu - W_0$①，其中，反向遏止电压 $U = \frac{E_k}{e}$②，结合题图可得 $v_a < v_c < v_b$，光的色散规律如图所示，从上往下依次为红、橙、黄、绿、青、蓝、紫，即红光的折射率最小，频率最小，紫光的频率、折射率最大，故 C 正确。
\onepicture[h=2.6cm]{figs/18_an.jpg}
}

\item
%% number: 19
%% paperName: 2019年天津理综第 19 题 6 分
%% typeId: 2
%% type: 多选题
%% chapter: 原子和原子核
%% point: 核聚变
%% method: 无
%% score: 6
%% degree: 750
%% duplicateId: 0
%% body:
我国核聚变反应研究大科学装置“人造太阳”2018 年获得重大突破，等离子体中心电子温度首次达到 1 亿度，为人类开发利用核聚变能源奠定了重要的技术基础。下列关于聚变的说法正确的是\xzanswer{AD}
\onepicture[h=4.2cm]{figs/19.jpg}
\fourchoices[answer=AD]
{核聚变比核裂变更为安全、清洁}
{任何两个原子核都可以发生聚变}
{两个轻核结合成质量较大的核，总质量较聚变前增加}
{两个轻核结合成质量较大的核，核子的比结合能增加}
%% answer: AD
%% memo:
\memoanswer{
重核裂变燃料及废物都需要严格处理，防止泄漏。轻核聚变原料及生成物都是安全的，更加清洁，故 A 正确；只有较轻的一些原子核可以实现聚变反应，故 B 错误；核聚变的过程中有质量亏损，比结合能增加，放出大量的能量，故 C 错误 D 正确。
}

\item
%% number: 20
%% paperName: 2019年天津理综第 20 题 6 分
%% typeId: 2
%% type: 多选题
%% chapter: 机械波
%% point: 波的图象
%% method: 无
%% score: 6
%% degree: 450
%% duplicateId: 0
%% body:
一列简谐横波沿 $x$ 轴传播，已知 $x$ 轴上 $x_{1} = 1\Um$ 和 $x_{2} = 7\Um$ 处质点的振动图像分别如图 1、图 2 所示，则此列波的传播速率可能是\xzanswer{BC}
\onepicture[h=3.6cm]{figs/20.png}
\fourchoices[answer=BC]
{$7\Ums$}
{$2\Ums$}
{$1.2\Ums$}
{$1\Ums$}
%% answer: BC
%% memo:
\memoanswer{
由题图 2 可知，波的周期 $T = 4\Us$。若波沿 $x$ 轴正方向传播，则由题图可知，$x = 7\Um$ 处质点振动滞后 $x = 1\Um$ 处质点的时间 $t = \left(n + \frac{1}{4}\right)T$（$n = 0,1,2,\ldots$），所以波传播的速度为 $v = \frac{x_2 - x_1}{t} = \frac{6\Um}{(4n+1)\Us}$（$n = 0,1,2,\ldots$），当 $n = 1$ 时，$v = 1.2\Ums$，故 C 正确；若波沿 $x$ 轴负方向传播，同理有 $t = \left(n + \frac{3}{4}\right)T$（$n = 0,1,2,\ldots$），所以波速 $v = \frac{6\Um}{(4n+3)\Us}$（$n = 0,1,2,\ldots$），当 $n = 0$ 时 $v = 2\Ums$，故 B 正确。
}

\item
%% number: 21
%% paperName: 2019年天津理综第 21 题 6 分
%% typeId: 2
%% type: 多选题
%% chapter: 电磁感应
%% point: 法拉第电磁感应定律
%% method: 无
%% score: 6
%% degree: 450
%% duplicateId: 0
%% body:
单匝闭合矩形线框电阻为 $R$，在匀强磁场中绕与磁感线垂直的轴匀速转动，穿过线框的磁通量 $\Phi$ 与时间 $t$ 的关系图像如图所示。下列说法正确的是\xzanswer{BC}
\onepicture[h=3.6cm]{figs/21.png}
\fourchoices[answer=BC]
{$\frac{T}{2}$ 时刻线框平面与中性面垂直}
{线框的感应电动势有效值为 $\frac{\sqrt{2}\pi\Phi_m}{T}$}
{线框转一周外力所做的功为 $\frac{2\pi^{2}\Phi_m^{2}}{RT}$}
{从 $t = 0$ 到 $t = \frac{T}{4}$ 过程中线框的平均感应电动势为 $\frac{\pi\Phi_m}{T}$}
%% answer: BC
%% memo:
\memoanswer{
由题图可知，$\frac{T}{2}$ 时刻磁通量反向最大，线框仍与中性面重合，故 A 错误；角速度 $\omega = \frac{2\pi}{T}$，最大感应电动势 $E_{\mathrm{m}} = n\Phi_{\mathrm{m}}\omega = \frac{2\pi\Phi_{\mathrm{m}}}{T}$，则感应电动势的有效值 $U = \frac{E_{\mathrm{m}}}{\sqrt{2}} = \frac{\sqrt{2}\pi\Phi_{\mathrm{m}}}{T}$，故 B 正确；线框匀速转动，外力做功转化为电能，求电能时所用电动势为有效值，所以 $W_{\text{外}} = \frac{U^2}{R}T = \frac{2\pi^2\Phi_{\mathrm{m}}^2}{RT}$，故 C 正确；$0 \sim \frac{T}{4}$ 过程中平均感应电动势 $\overline{E} = n\frac{\Delta\Phi}{\Delta t} = \frac{\Phi_{\mathrm{m}}}{\frac{T}{4}} = \frac{4\Phi_{\mathrm{m}}}{T}$，故 D 错误。
}

\item
%% number: 22
%% paperName: 2019年天津理综第 22 题 6 分
%% typeId: 4
%% type: 实验题
%% chapter: 电路
%% point: 电阻定律
%% method: 无
%% score: 6
%% degree: 550
%% duplicateId: 0
%% body:
\begin{enumerate}
	\item
	第 26 届国际计量大会决定，质量单位“千克”用普朗克常量 $h$ 定义，“国际千克原器”于 2019 年 5 月 20 日正式“退役”。$h$ 的数值为 $6.63\times10^{-34}$，根据能量子定义，$h$ 的单位是\tkanswer[1.8cm]{$\UJcdots$}，该单位用国际单位制中的力学基本单位表示，则为\tkanswer[2.4cm]{$\Ukgmqs$}。
	\item
	某小组做测定玻璃的折射率实验，所用器材有：玻璃砖，大头针，刻度尺，圆规，笔，白纸。

	\begin{enumerate}
		\item
		下列哪些措施能够提高实验准确程度\tkanswer{AD}。

		（A）选用两光学表面间距大的玻璃砖

		（B）选用两光学表面平行的玻璃砖

		（C）选用粗的大头针完成实验

		（D）插在玻璃砖同侧的两枚大头针间的距离尽量大些
		\item
		该小组用同一套器材完成了四次实验，记录的玻璃砖界线和四个大头针扎下的孔洞如下图所示，其中实验操作正确的是\tkanswer{D}。

		\fourchoices[answer=D, ispicture=true, h=2.2cm]
		{\includesvg{figs/22a}}
		{\includesvg{figs/22b}}
		{\includesvg{figs/22c}}
		{\includesvg{figs/22d}}
		\item
		该小组选取了操作正确的实验记录，在白纸上画出光线的径迹，以入射点 O 为圆心作圆，与入射光线、折射光线分别交于 A、B 点，再过 A、B 点作法线 NN$'$ 的垂线，垂足分别为 C、D 点，如图所示，则玻璃的折射率 $n=$\tkanswer[1.8cm]{$\frac{AC}{BD}$}。（用图中线段的字母表示）

		\onepicture[h=4.0cm]{figs/22e.png}
	\end{enumerate}
	\item
	现测定长金属丝的电阻率。

	\begin{enumerate}
		\item
		某次用螺旋测微器测量金属丝直径的结果如图所示，其读数是\tkanswer[1.6cm]{0.200}$\Umm$。

		\onepicture[align=right, h=2.4cm]{figs/22f.png}
		\item
		利用下列器材设计一个电路，尽量准确地测量一段金属丝的电阻。这段金属丝的电阻 $R_{x}$，约为 $100\UO$，画出实验电路图，并标明器材代号。

		电源 $E$\quad（电动势 $10\UV$，内阻约为 $10\UO$）

		电流表 A$_{1}$\quad（量程 $0 \sim 250\UmA$，内阻 $R_{1} = 5\UO$）

		电流表 A$_{2}$\quad（量程 $0 \sim 300\UmA$，内阻约为 $5\UO$）

		滑动变阻器 $R$\quad（最大阻值 $10\UO$，额定电流 $2\UA$）

		开关 S 及导线若干
		\item
		某同学设计方案正确，测量得到电流表 A$_{1}$ 的读数为 $I_{1}$，电流表 A$_{2}$ 的读数为 $I_{2}$，则这段金属丝电阻的计算式 $R_{x}=$\tkanswer[2.0cm]{$\frac{I_1R_1}{I_2-I_1}$}。从设计原理看，其测量值与真实值相比\tkanswer[1.4cm]{相等}（填“偏大”、“偏小”或“相等”）。
	\end{enumerate}
\end{enumerate}
%% answer:
\jdanswer{
\begin{enumerate}
	\item
	$\UJcdots$，$\Ukgmqs$
	\item
	① AD；② D；③ $\frac{AC}{BD}$
	\item
	① $0.200\Umm$（$0.196 \sim 0.204\Umm$ 均可）；② 如图所示；③ $\frac{I_1R_1}{I_2-I_1}$，相等
\end{enumerate}
}
%% memo:
\memoanswer{
（1）由能量子公式 $E = h\nu$，能量的单位为 $\UJ$，频率的单位为 $\Usn$，故 $h$ 的单位为 $\UJcdots$；又因能量的单位换成力学基本单位可表示为 $\Ukgmqsq$，则 $h$ 的单位为 $\Ukgmqs$。

（2）①玻璃砖越厚，入射光线与出射光线侧移距离越大，作图时相对误差越小；两枚大头针间的距离越大，在确定入射及出射光线时的相对误差越小；大头针应选用较细的，减少测量误差；玻璃砖两光学面平行与否不影响测量结果。故 AD 正确。

②出射光线应向左侧移且平行于入射光线，故 D 正确。

③设入射角为 $\theta_1$，折射角为 $\theta_2$，则折射率 $n=\frac{\sin\theta_1}{\sin\theta_2}=\frac{\frac{AC}{AO}}{\frac{BD}{BO}}=\frac{AC}{BD}$。

（3）①由题图可知，螺旋测微器的读数为 $0\Umm + 20.0 \times 0.01\Umm = 0.200\Umm$。

②实验电路图如答图所示。

\onepicture[h=3.0cm]{figs/22_an.png}

③电流表 A$_{1}$ 与金属丝 $R_{x}$ 并联，则 $R_{x}$ 两端电压为 $I_{1}R_{1}$，通过 $R_{x}$ 的电流为 $I_{2}-I_{1}$，则金属丝的阻值 $R_{x}=\frac{I_{1}R_{1}}{I_{2}-I_{1}}$。根据分析可知，理论上测量值与真实值相等。
}

\item
%% number: 23
%% paperName: 2019年天津理综第 23 题 16 分
%% typeId: 6
%% type: 计算题
%% chapter: 功和能
%% point: 动能定理
%% method: 无
%% score: 16
%% degree: 650
%% duplicateId: 0
%% body:
完全由我国自行设计、建造的国产新型航空母舰已完成多次海试，并取得成功。航母上的舰载机采用滑跃式起飞，故甲板是由水平甲板和上翘甲板两部分构成，如图 1 所示。为了便于研究舰载机的起飞过程，假设上翘甲板 BC 是与水平甲板 AB 相切的一段圆弧，示意如图 2，AB 长 $L_{1}=150\Um$，BC 水平投影 $L_{2} = 63\Um$，图中 C 点切线方向与水平方向的夹角 $\theta = 12^\circ$（$\sin 12^\circ \approx 0.21$）。若舰载机从 A 点由静止开始做匀加速直线运动，经 $t = 6\Us$ 到达 B 点进入 BC。已知飞行员的质量 $m = 60\Ukg$，$g = 10\Umsq$，求：
\begin{enumerate}
	\item
	舰载机水平运动的过程中，飞行员受到的水平力所做功 $W$；
	\item
	舰载机刚进入 BC 时，飞行员受到竖直向上的压力 $F_{N}$ 多大。
\end{enumerate}
\onepicture[h=4.2cm, align=right]{\includesvg{figs/23}}
%% answer:
\jdanswer{
\begin{enumerate}
	\item
	$W = 7.5\times10^{4}\UJ$
	\item
	$F_{N} = 1.1\times10^{3}\UN$
\end{enumerate}
}
%% memo:
\memoanswer{
（1）舰载机由静止开始做匀加速直线运动，设其刚进入上翘甲板时的速度为 $v$，则有

$\frac{v}{2}=\frac{L_{1}}{t}$①，

由动能定理得

$W=\frac{1}{2}mv^{2}-0$②，

联立①②式，代入数据，得

$W = 7.5 \times 10^{4}\UJ$③。

（2）设上翘甲板所对应的圆弧半径为 $R$，由几何关系得

$L_{2}=R\sin\theta$④，

由牛顿第二定律得

$F_{N}-mg=m\frac{v^{2}}{R}$⑤，

联立①④⑤式，代入数据得

$F_{N}=1.1\times10^{3}\UN$⑥。
}

\item
%% number: 24
%% paperName: 2019年天津理综第 24 题 18 分
%% typeId: 6
%% type: 计算题
%% chapter: 电磁感应
%% point: 电磁感应能量分析
%% method: 无
%% score: 18
%% degree: 450
%% duplicateId: 0
%% body:
如图所示，固定在水平面上间距为 $l$ 的两条平行光滑金属导轨，垂直于导轨放置的两根金属棒 MN 和 PQ 长度也为 $l$、电阻均为 $R$，两棒与导轨始终接触良好。MN 两端通过开关 S 与电阻为 $R$ 的单匝金属线圈相连，线圈内存在竖直向下均匀增加的磁场，磁通量变化率为常量 $k$。图中虚线右侧有垂直于导轨平面向下的匀强磁场，磁感应强度大小为 $B$。PQ 的质量为 $m$，金属导轨足够长，电阻忽略不计。
\begin{enumerate}
	\item
	闭合 S，若使 PQ 保持静止，需在其上加多大的水平恒力 $F$，并指出其方向；
	\item
	断开 S，PQ 在上述恒力作用下，由静止开始到速度大小为 $v$ 的加速过程中流过 PQ 的电荷量为 $q$，求该过程安培力做的功 $W$。
\end{enumerate}
\onepicture[h=3.0cm, align=right]{figs/24.png}
%% answer:
\jdanswer{
\begin{enumerate}
	\item
	$F = \frac{Bkl}{3R}$，方向水平向右
	\item
	$W = \frac{1}{2}mv^{2} - \frac{2}{3}kq$
\end{enumerate}
}
%% memo:
\memoanswer{
（1）设线圈中的感应电动势为 $E$，由法拉第电磁感应定律 $E=\frac{\Delta\Phi}{\Delta t}$，则 $E=k$①，

设 PQ 与 MN 并联的电阻为 $R_{\text{并}}$，有

$R_{\text{并}}=\frac{R}{2}$②，

闭合 S 时，设线圈中的电流为 $I$，由闭合电路欧姆定律得

$I=\frac{E}{R_{\text{并}}+R}$③，

设 PQ 中的电流为 $I_{PQ}$，有

$I_{PQ}=\frac{1}{2}I$④，

设 PQ 受到的安培力为 $F_{\text{安}}$，有

$F_{\text{安}}=BI_{PQ}l$⑤，

保持 PQ 静止，由平衡条件得

$F=F_{\text{安}}$⑥，

联立①②③④⑤⑥式得

$F=\frac{Bkl}{3R}$⑦，方向水平向右。

（2）设 PQ 由静止开始到速度大小为 $v$ 的加速过程中，PQ 运动的位移为 $x$，所用时间为 $\Delta t$，回路中的磁通量变化为 $\Delta\Phi$，平均感应电动势为 $\overline{E}=\frac{\Delta\Phi}{\Delta t}$⑧，其中 $\Delta\Phi=Blx$⑨，

设 PQ 中的平均电流为 $\overline{I}=\frac{\overline{E}}{2R}$⑩，

由电流的定义得 $\overline{I}=\frac{q}{\Delta t}$\ccnum{⑪}，

由动能定理得 $Fx + W = \frac{1}{2}mv^2 - 0$\ccnum{⑫}，

联立⑦⑧⑨⑩\ccnum{⑪⑫}式得 $W = \frac{1}{2}mv^2 - \frac{2}{3}kq$\ccnum{⑬}。
}

\item
%% number: 25
%% paperName: 2019年天津理综第 25 题 20 分
%% typeId: 6
%% type: 计算题
%% chapter: 运动和力
%% point: 牛顿定律的应用
%% method: 无
%% score: 20
%% degree: 450
%% duplicateId: 0
%% body:
2018 年，人类历史上第一架由离子引擎推动的飞机诞生，这种引擎不需要燃料，也无污染物排放。引擎获得推力的原理如图所示，进入电离室的气体被电离成正离子，而后飘入电极 A、B 之间的匀强电场（初速度忽略不计），A、B 间电压为 $U$，使正离子加速形成离子束，在加速过程中引擎获得恒定的推力。单位时间内飘入的正离子数目为定值，离子质量为 $m$，电荷量为 $Ze$，其中 $Z$ 是正整数，$e$ 是元电荷。
\begin{enumerate}
	\item
	若引擎获得的推力为 $F_{1}$，求单位时间内飘入 A、B 间的正离子数目 $N$ 为多少；
	\item
	加速正离子束所消耗的功率 $P$ 不同时，引擎获得的推力 $F$ 也不同，试推导 $\frac{F}{P}$ 的表达式；
	\item
	为提高能量的转换效率，要使 $\frac{F}{P}$ 尽量大，请提出增大的三条建议。
\end{enumerate}
\onepicture[h=3.0cm, align=right]{\includesvg{figs/25}}
%% answer:
\jdanswer{
\begin{enumerate}
	\item
	$N = \frac{F_{1}}{\sqrt{2ZemU}}$
	\item
	$\frac{F}{P} = \sqrt{\frac{2m}{ZeU}}$
	\item
	用质量大的离子；用带电量少的离子；减小加速电压。
\end{enumerate}
}
%% memo:
\memoanswer{
（1）设正离子经过电极 B 时的速度为 $v$，由动能定理得

$ZeU = \frac{1}{2}mv^{2} - 0$①，

设正离子束所受的电场力为 $F_{1}^{\prime}$，由牛顿第三定律得

$F_{1}^{\prime}=F_{1}$②，

设引擎在 $\Delta t$ 时间内飘入电极间的正离子个数为 $\Delta N$，由牛顿第二定律得

$F_{1}^{\prime}=\Delta Nm\frac{v-0}{\Delta t}$③，

联立①②③式，且 $N=\frac{\Delta N}{\Delta t}$ 得

$N=\frac{F_{1}}{\sqrt{2ZemU}}$④。

（2）设正离子束所受的电场力为 $F'$，由正离子束在电场中做匀加速直线运动，有

$P = \frac{1}{2} F' v$⑤，

考虑到牛顿第三定律得到 $F' = F$，联立①⑤式得

$\frac{F}{P} = \sqrt{\frac{2m}{ZeU}}$⑥。

（3）为使 $\frac{F}{P}$ 尽量大，分析⑥式得到三条建议：用质量大的离子；用带电量少的离子；减小加速电压。
}

\end{enumerate}

\end{document}
